% ===== 이번 주차 정리 =====
\begin{frame}{이번 주차 정리}
  \small
  \begin{enumerate}
    \item \kb{1.1 역사와 로드맵}
      \begin{itemize}
        \item 역사는 직선이 아님 --- \textbf{두 번의 AI 겨울}(1969 XOR
          증명, 1986 역전파로 해소)과 2012 AlexNet.
        \item 돌파구의 반복 패턴: \textbf{새 수학이 아니라 데이터+
          계산의 재조합}. Block A(Week 2$\sim$7 고전 ML) $\to$
          Block B(Week 9$\sim$16 신경망$\to$생성모델).
      \end{itemize}
    \item \kb{1.2 문제 정식화와 세 갈래}
      \begin{itemize}
        \item 알고리즘보다 먼저: \textbf{$x$ · $y$ · $h$ 로 정의}.
          ``라벨이 있는가? 보상이 있는가?'' 두 질문으로
          \textbf{지도/비지도/강화}를 가름.
        \item 회귀 vs 분류 = \textbf{같은 모델, 다른 출력}(항등 vs
          시그모이드 링크) · 모든 알고리즘의 뼈대 =
          \textbf{$J$ 정의 $\to$ 최소화}.
      \end{itemize}
    \item \kb{1.3 파이프라인과 사전지식}
      \begin{itemize}
        \item 모델링은 5단계 중 하나 --- 더티 데이터(결측·중복·
          이상치)·EDA·\textbf{train/test 순서(누수 방지)}가 모델의
          상한을 결정.
        \item 수학 전량 = \textbf{내적 · 연쇄법칙 · 베이즈 정리}.
      \end{itemize}
  \end{enumerate}
  \vskip0.6em
  \begin{block}{다음 주차 --- Week 2. 회귀 모델: 선형회귀와
  로지스틱회귀}
    이 책의 첫 번째 모델. 그 도입부에는 이미 19세기 수학자가 잡음
    섞인 점들 사이를 ``모두를 가장 잘 맞혀주는'' 한 줄로 통과시켜야
    할 문제가 등장한다 --- 오늘 세운 ``\textbf{손실함수를 정의하고
    최소화한다}''는 틀과 \textbf{내적·그래디언트 감각}이 있으면,
    그 문제는 낯선 수학이 아니라 바로 이 주차의 \textbf{연장선}으로
    읽힌다.
  \end{block}
\end{frame}
